\section{Coprime Pell fusion: ten positive variables and degree 266}
\label{sec:positive-fusion}

The following construction improves the pinned ten-positive-variable
endpoints from $(103,340)$ and $(98,368)$, where the ordered pair records
live binary arithmetic operations and exact total degree, to
\[
                    \boxed{(96,266)}.
\]
The operation split is 57 multiplications, 20 additions and 19 subtractions.
Multiplication by a constant counts as an operation. This is an ordinary
unbounded arithmetic reduction relative to the retained source/native
calibration, not a new Lean or encoded internal-ZFC certificate. No
consistency or global minimality claim is made.

\subsection{Complete polynomial}
All ten variables
\[
       J,F,\alpha,i,T,x,Y,n,\eta,\zeta
\]
range over strictly positive integers. The four fixed integer coefficients
are the complete original compact44 tuple bound in \texttt{equation.json}
and its coefficient payloads. In particular,
\[
 \beta\ge31\text{ is odd},\qquad K_0>0,\qquad
 0<M_C<\beta<M_F<2\beta.
\]
Here $M_F$ already includes its source shift. The distinct repaired BFS44
tuple must not be mixed with this tuple; its $K_0$ is not the same integer.
Define the following polynomial abbreviations:
\begin{align*}
q&=\beta J+1,& X&=q+x,& \mu&=q^2+1,\\
A&=\mu Y(X+1)+1,&\Delta&=A^2-1,&H&=4A-5,\\
k&=\eta+\zeta,&c&=kY+\eta,&V_0&=X\mu^2Y^2,\\
D&=V_0(V_0+1)k^2+1,&G&=(q-1)(q-F)+\alpha,\\
R&=G(q^2-1)+(M_C+qM_F)J,&e&=k-R-1,\\
C&=q-F-\alpha,&B&=(qK_0+X)C-qF,\\
E&=ic^2,&\mathcal F&=1+\Delta DE^2,&Q&=\Delta^2DE^2,\\
\delta&=Q-1,&b&=c(T+\alpha)-R,&V&=\mathcal F b-c,\\
D_y&=QV^2-1,&L&=HDE+b,&W&=1+q(q-1)L,\\
P&=q^4(q-1),&M&=P\mu^2H,&K_+&=X(3c-2X)+2.
\end{align*}
Finally set
\begin{align*}
Z={}&P\mu^2\{H(n+V)+K_+\}\\
 &+H\{\mu^2[q(q-1)Y-q^3WB-(q-1)X]-Pe\},\\
a={}&D\mathcal F,\qquad s=\delta Z^2,\qquad
u=\delta M^2a,\qquad v=M^2D_y.
\end{align*}
The single equation is
\begin{equation}\label{eq:fused-positive}
             \boxed{\mathcal N=(s-v-u)^2-4uv=0.}
\end{equation}
There are no additional quantified abbreviation variables and no radicals
or divisions in its evaluator. The complete circuit uses the stored names
\texttt{Dtau}$=D$, \texttt{Df}$=\mathcal F$, and \texttt{I}$=E$.
The latter is \emph{not} the actual strong Pell ordinate introduced in the
proof below.

The structural change is to make the first radicand $D$ a factor of the
strong squared ordinate. Since $\mathcal F\equiv1\pmod D$, a square
product $D\mathcal F$ forces its two factors to be squares separately.
This permits a two-radical quartic norm in place of the parent's
three-radical octic norm.

\subsection{Exact contraction of the literal evaluator}
The first fused evaluator had 98 operations; the current 96-operation
schedule is identically equal to it. Write $U=\mu^2$, $S=q^3$,
$g=q(q-1)$ and $D_{\rm low}=(q-1)X$, so $P=Sg$ and $W=1+gL$.
The literal evaluator uses
\[
 Z=PU\{H(n+V-LB)+K_+\}
       +H\{U(gY-SB-D_{\rm low})-Pe\}.
\]
Subtracting the earlier displayed expression gives
$HULB(Sg-P)=0$ universally. Delete the private $gL,W,WB$ definitions
and insert only $LB$ and $(n+V)-LB$. Also compute
$P=q^3[q(q-1)]$ using the existing $q(q-1)$ register, deleting the
private $q^4$ register. This saves two operations; all remaining kernel
and norm gates are unchanged. The final split is 57 multiplications,
20 additions and 19 subtractions. The abbreviation $W$ remains useful
in the proof but is not separately evaluated. The complete polynomial,
its degree, and both positive witness maps are unchanged. This same
contraction appears in upstream commit 49379c31d, yielding comparison
endpoints 98/368 and 103/340; the fused endpoint is 96/266.

\subsection{Root extraction and the implicit strong ordinate}
The exact rational carrier identity is
\begin{equation}\label{eq:fused-center}
\frac ZM=n+V+\frac{K_+}{H}+\frac Y{q^3}
 -\frac{WB}{q(q-1)}-\frac X{q^4}-\frac e{\mu^2}.
\end{equation}
For every positive input, $q\ge32$, $X>q$, $A>4$, $c\ge3$, $E\ge c^2$
and $D>1$. The parent's elementary adjacent-radicand bounds, strengthened
by the new factor $D$, give
\[
 \delta>c^2-1>0,\qquad \mathcal F>c+2,\qquad \mathcal F<\delta.
\]
If $b=0$, then $V=-c$ and $c^2<D_y/\delta<c^2+1$. If $b\ne0$, then
$|V|\ge\mathcal F-c>1$. Thus $a$ and $D_y/\delta$ are positive
rational numbers before integrality is established.

They are unequal: $\delta a=D_y$ would reduce modulo $D>1$ to
$0=-1$, since $D\mid Q$ and $D\mid a$. The elementary norm identity
\[
 (z^2-r_1^2-r_2^2)^2-4r_1^2r_2^2
 =\prod_{\epsilon,\nu\in\{-1,1\}}(z-\epsilon r_1-\nu r_2)
\]
applied to \eqref{eq:fused-positive} therefore gives a nonzero rational
signed sum
\[
 \frac ZM=\epsilon\sqrt a+\nu\sqrt{D_y/\delta}.
\]
Both radicals are rational: squaring this relation and solving linearly
for either radical proves the assertion, since $Z/M\ne0$. As $a$ is an
integer, it is an integer square. Furthermore
$\gcd(D,\mathcal F)=1$, so unique factorization supplies positive integers
\[
       D=\tau^2,\qquad \mathcal F=f^2.
\]
Define the actual strong ordinate
\[
       I_* = \tau E.
\]
It is an integer even though it is not asserted to be a polynomial in the
ten inputs. Then
\[
 f^2-\Delta I_*^2=1,\qquad Q=(\Delta I_*)^2,\qquad
 c^2\mid I_*,\qquad I_*\ge c.
\]
The retained adjacent-coefficient rational-to-integral descent applied to
$QV^2-(Q-1)y^2=1$ proves that $y=\sqrt{D_y/\delta}$ is a positive integer.
Consequently the carrier relation is
\begin{equation}\label{eq:fused-signed}
                  Z/M=\epsilon\tau f+\nu y.
\end{equation}
None of these steps assumes parity of $q$, transport positivity, or
$\mu^2\mid e$.

\subsection{Soundness with delayed denominator cancellation}
Write $\chi_A(t),\psi_A(t)$ for the standard Pell sequences defined by
\[
 \chi_A(t)+\psi_A(t)\sqrt{A^2-1}
       =(A+\sqrt{A^2-1})^t.
\]
The retained all-index strong/auxiliary folding lemma applies to
$A,c,I_*,f,V,y$: its hypotheses include $A\ge3$, $c\le I_*$,
$f>2I_*$ and $V=-c\pmod f$. It gives
$c=\psi_A(p)$ for an odd $p\ge1$. The first Pell equation gives
\[
 k=2\psi_{2V_0+1}(n_0),\qquad
 \tau=\chi_{2V_0+1}(n_0),\qquad n_0\ge1.
\]
The case $n_0=1$ would give $c=2Y+1$ strictly between
$\psi_A(1)=1$ and $\psi_A(2)=2A$. Hence
\begin{equation}\label{eq:fused-cbound}
 c>kY\ge k\ge4(2V_0+1)>8X\mu^2Y^2.
\end{equation}
The unchanged mask estimate gives $R<c\alpha/4$. As in the parent this
implies
\[
 b>0,\quad V>0,\quad |e/\mu^2|<b/2,\quad 0<y-V<b/2.
\]
The adjacent-root estimate uses $\mathcal F<\delta$. Also $c>2X$, so
$K_+>0$.

The new polynomial transport majorant is adequate because $f<AI_*$
and therefore $\tau f<ADE$. Since $D>X$ and $E\ge1$,
\begin{equation}\label{eq:fused-majorant}
       \tau f+X/q^4<(A+1)DE<HDE.
\end{equation}
If $B\le-1$, then $WB/[q(q-1)]<-L$. Solving
\eqref{eq:fused-center}--\eqref{eq:fused-signed} for $n$ and using
\eqref{eq:fused-majorant}, the two $b/2$ bounds, $K_+/H>0$ and $Y/q^3>0$
gives $n<HDE+b-L=0$, a contradiction. Thus $B\ge0$, and its definition
forces $C>0$. The unchanged mask bounds now give
\[
                         3q<R<q^4-2.
\]
The normalized strong-rank and signed auxiliary-index lemmas apply since
$c^2\mid I_*$, $c>2R$, $c>2p$, $V=-c\pmod f$ and $V=-R\pmod c$.
They give $p=R$. Oddness of $p$ gives $R$ odd, and the original mask
parity calculation gives $q$ even and $J$ odd.

Only at this point use
\[
                   \gcd(\mu^2,q^4(q-1)H)=1.
\]
Indeed $H=-1\pmod\mu$, $\gcd(q,\mu)=1$, and evenness of $q$ handles
$q-1$. Since \eqref{eq:fused-signed} is integral, separating the single
modulus $\mu^2$ in \eqref{eq:fused-center} yields $\mu^2\mid e$.
The first Pell congruence and the unchanged strict comparison
$2V_0+1>A$ give $n_0=(R+1)/2$. Put $R=2r+1$; then $r\ge48$.

The retained zero-offset ratio lemma, uniform in $X\ge32$, is
\[
 \xi<\frac{\psi_A(2r+1)}{\psi_{2X\mu^2Y^2+1}(r+1)}
       <\xi\left(1+\frac1{Y(X+1)}\right)^{2r},
 \qquad \xi=\frac{(X+1)^{2r}}{X^r}.
\]
It gives $Y>X^r/3$ and $X^{r+1}<3Y(X+1)<H$. From the integral center,
$H\mid PK_+$. Since $H$ is odd and
\[
 K_+=HXc-2K_2,\qquad K_2=X[X+2(A-2)c]-1,
\]
cancel only 2 and apply the universal projection lemma:
\[
          H\mid q^4(q-1)(X-2^R)(2^RX+1).
\]
The whole product has absolute value less than $X^{r+1}<H$ when
$X\ge q\ge32$, $r\ge48$, as in the direct-base parent. Thus $X=2^R$.
This is whole-product vanishing, not cancellation of $q$ through $H$ or
a choice of a square-root branch separately at each prime. Projection
now also makes $K_+/H$ integral.

Removing the integral terms from \eqref{eq:fused-center} leaves
$Y/q^3-WB/[q(q-1)]-X/q^4$ integral. Clear the denominator and reduce
modulo $q$ to obtain $q\mid X$. Hence $q=2^\kappa$; since
$R>3q>4\kappa$, also $q^4\mid X$. Set
\[
 w=X/q,\qquad B_t=(K_0+w)C-F,\qquad B=qB_t.
\]
The remaining partial fractions give $q^3\mid Y$ and $q-1\mid B_t$,
using $W=1\pmod{q-1}$. Thus $t=B_t/(q-1)\ge0$ is an integer.
The fixed-minus auxiliary parity lemma gives $R=3\pmod4$, and the
same ratio rounding recovers the source half-binomial $Y$.

The parent's fresh unscaled source reconstruction, at
$A_{\mathrm{source}}=Y(X+1)+2$, now applies to the recovered data.
It does not reuse $I_*$ or equate the two Pell parameters. The pinned
source61/native/controller calibration yields a finite syntactic
refutation. This proves soundness with exactly that inherited calibration.

\subsection{Completeness by a new strong Pell period}
A finite refutation supplies the parent's canonical source data with
$X=2^R>q$, $R=3\pmod4$, $r\ge48$, $q^3\mid Y$, $q^4\mid X$ and
$t\ge0$. Retain $J,F,\alpha,Y$ and set $x=X-q>0$. The unchanged
zero-offset ratio lemma supplies positive integers
\[
 c=\psi_A(R),\quad k=2\psi_{2V_0+1}(r+1),\quad
 \tau=\chi_{2V_0+1}(r+1),
\]
\[
 \eta=c-kY>0,\qquad\zeta=k-\eta>0,\qquad
 \rho=(k-R-1)/\mu^2\ge1.
\]
Choose a positive $m$ such that $\tau c^2\mid\psi_A(m)$. This exists
by finiteness: the determinant-one recurrence matrix
\[
          \begin{pmatrix}2A&-1\\1&0\end{pmatrix}
\]
has a positive power equal to the identity modulo $\tau c^2$. Its
corresponding Pell ordinate is zero modulo that modulus. Set
\[
 I_*=\psi_A(m),\qquad i=\frac{I_*}{\tau c^2}>0,
          \qquad f=\chi_A(m).
\]
Then $E=ic^2$, $D=\tau^2$ and
$\mathcal F=1+\Delta DE^2=f^2$. No divisibility of the parent's
former strong ordinate is assumed.

The retained universal auxiliary-completeness lemma applies to this
arbitrary new strong pair because $R=3\pmod4$, $c=\psi_A(R)$ and
$c\mid I_*$. It supplies positive $T_*,u_*$ with
\[
 V=f^2(cT_*-R)-c,\qquad y=V+u_*,\qquad QV^2-(Q-1)y^2=1.
\]
The same explicit growth estimate gives $T_*>\alpha$; put
$T=T_*-\alpha>0$. Thus the auxiliary witness is also genuinely rebuilt.

Projection gives an integer $k_+=K_+/H$ with $0<k_+<Xc$.
As in the parent, $f>(A-1)I_*>Xc$ and $\tau>Y\ge Y/q^3$.
Since $\tau,f\ge3$, $\tau f>\tau+f>Y/q^3+k_+$. Define
\[
 n=Wt+\rho+\tau f+u_*-Y/q^3-k_++X/q^4.
\]
It is a positive integer, including when $t=0$. Substitution into
\eqref{eq:fused-center} gives $Z/M=\tau f+y$, so one factor of the
quartic norm vanishes. This constructs all ten positive coordinates
from every refutation. A pointwise positive inverse between the old
and new witness coordinates is neither needed nor asserted.

\subsection{Exact degree and reproducible scope}
Give each variable degree one and each fixed coefficient degree zero.
The degree bounds for the polynomial abbreviations are
\[
 A:4,\ D:16,\ E:5,\ \mathcal F:34,\ Q:42,\ V:38,\ L:25,
 \ W:27,\ M:13,\ Z:51.
\]
The uncancelled syntactic bound on $\mathcal N$ is 288. Put $Z_0=Z-MV$;
its degree is 40, with highest part furnished uniquely by
$-H\mu^2q^3WB$. The exact identity
\[
 \omega:=s-v
  =\delta Z_0^2+2\delta MVZ_0-M^2(V^2-1)
\]
has summand bounds 122, 133 and 102. Also $\deg u=118$ and
$\deg v=144$. Hence
\[
 \mathcal N=(\omega-u)^2-4uv
\]
has unique top candidate $\omega^2$ of degree266; the other term bounds
are 251,236 and262.

Set $q_0=\beta J$, $X_0=q_0+x$, $k_0=\eta+\zeta$ and
\[
 B_0=(q_0K_0+X_0)(q_0-F-\alpha)-q_0F.
\]
The complete leading form is
\[
 \omega_0=2^7q_0^{63}X_0^{15}i^5Y^{31}k_0^{16}B_0(q_0-F).
\]
The highest power of $x$ in $\omega_0^2$ is32. Selecting
$\alpha^2F^2\eta^{32}$ in the remaining factors gives
\[
 [J^{126}x^{32}i^{10}Y^{62}\eta^{32}\alpha^2F^2],\mathcal N
                  =2^{14}\beta^{126}>0.
\]
The exponents sum to 266, proving the exact total degree at the actual
fixed tuple.

The portable packet contains the full circuit, coherent coefficient
payloads, copied ordinary dependencies, \texttt{build.py}, and a Python
standard-library-only independent \texttt{check.py}. Both normal and
optimized Python runs check the complete literal circuit, all fixed
payload hashes, the full norm and centering identities, 24 exact
independent-formula controls, actual-payload modular controls, nonzero
leading homogeneous witnesses, six finite Pell-period/fused-kernel cases,
64 signed whole-polynomial old/new comparisons, and five rejected mutations. These finite checks are transcription and
algebra checks; they do not replace the unbounded proof or the inherited
native/source calibration. The inherited arithmetic calibration is pinned to
repository commit \texttt{84e16e43092614ba}\allowbreak\texttt{e262ca59651cb9d77e31a255}.
